\documentclass{article}
\usepackage{mathtools} 
\usepackage{fontspec}
\usepackage[UTF8]{ctex}
\usepackage{amsthm}
\usepackage{mdframed}
\usepackage{xcolor}
\usepackage{amssymb}
\usepackage{amsmath}


% 定义新的带灰色背景的说明环境 zremark
\newmdtheoremenv[
  backgroundcolor=gray!10,
  % 边框与背景一致，边框线会消失
  linecolor=gray!10
]{zremark}{说明}

% 通用矩阵命令: \flexmatrix{矩阵名}{元素符号}{行数}{列数}
\newcommand{\flexmatrix}[4]{
  \[
  #1 = \begin{pmatrix}
    #2_{11}     & #2_{12}     & \cdots & #2_{1#4}   \\
    #2_{21}     & #2_{22}     & \cdots & #2_{2#4}   \\
    \vdots      & \vdots      & \ddots & \vdots     \\
    #2_{#31}    & #2_{#32}    & \cdots & #2_{#3#4}
  \end{pmatrix}
  \]
}

% 简化版命令（默认矩阵名为A，元素符号为a）: \quickmatrix{行数}{列数}
\newcommand{\quickmatrix}[2]{\flexmatrix{A}{a}{#1}{#2}}

\begin{document}
\title{习题 9.2}
\author{张志聪}
\maketitle

\section*{1}

\begin{itemize}
  \item (7)
        \begin{align*}
          \int_{0}^{4} \frac{dx}{1 + \sqrt{x}}
        \end{align*}
        令$t = \sqrt{x}, dx = 2t dt$，于是
        \begin{align*}
          \int \frac{dx}{1 + \sqrt{x}}
           & = \int \frac{2t}{1 + t} dt                               \\
           & = \int \frac{2t + 2 - 2}{1 + t} dt                       \\
           & = \int \frac{2t + 2}{1 + t} dt - \int \frac{2}{1 + t} dt \\
           & = 2t - 2 \ln |1 + t| + C                                 \\
           & = 2\sqrt{x} - 2 \ln |1 + \sqrt{x}| + C
        \end{align*}
        所以
        \begin{align*}
          \int_{0}^{4} \frac{dx}{1 + \sqrt{x}}
           & = (2\sqrt{x} - 2 \ln |1 + \sqrt{x}| + C) \bigg|_0^4 \\
           & = 4 - 2 \ln 3
        \end{align*}

\end{itemize}

\section*{2}

\begin{itemize}
  \item (1)

        我们有
        \begin{align*}
          \lim \limits_{n \to \infty} \frac{1}{n^4}(1 + 2^3 + \cdots + n^3)
           & = \lim \limits_{n \to \infty} \frac{1}{n} \left[(\frac{1}{n})^3 + (\frac{2}{n})^3 + \cdots + (\frac{n}{n})^3 \right] \\
           & = \lim \limits_{n \to \infty} \sum \limits_{i = 1}^n \frac{1}{n} \cdot \left(\frac{i}{n}\right)^3
        \end{align*}
        这个和式是函数$f(x) = x^3$在区间$[0, 1]$上的一个积分和。所以
        \begin{align*}
          \lim \limits_{n \to \infty} \frac{1}{n^4}(1 + 2^3 + \cdots + n^3)
           & = \int_{0}^{1} x^3 dx       \\
           & = \frac{1}{4}x^4 \bigg|_0^1 \\
           & = \frac{1}{4}
        \end{align*}

  \item (2)
        \begin{align*}
           & \lim \limits_{n \to \infty} n \left[ \frac{1}{(n + 1)^2} + \frac{1}{(n+2)^2} + \cdots + \frac{1}{(n+n)^2} \right]                  \\
           & = \lim \limits_{n \to \infty} \frac{1}{n}\left[ \frac{n^2}{(n + 1)^2} + \frac{n^2}{(n+2)^2} + \cdots + \frac{n^2}{(n+n)^2} \right] \\
           & = \lim \limits_{n \to \infty} \frac{1}{n}\left[ \left(\frac{n}{n + 1}\right)^2 + \left(\frac{n}{n+2}\right)^2
          + \cdots + \left(\frac{n}{n+n}\right)^2 \right]                                                                                       \\
           & = \lim \limits_{n \to \infty} \frac{1}{n}\left[ \left(\frac{1}{1 + \frac{1}{n}}\right)^2 + \left(\frac{1}{1+\frac{2}{n}}\right)^2
          + \cdots + \left(\frac{1}{1+\frac{n}{n}}\right)^2 \right]                                                                             \\
           & =\lim \limits_{n \to \infty} \sum \limits_{i=1}^n \frac{1}{n} \left(\frac{1}{1+\frac{i}{n}}\right)^2
        \end{align*}
        这个和式是函数$f(x) = \left(\frac{1}{1 + x}\right)^2$在区间$[0, 1]$上的一个积分和。所以
        \begin{align*}
           & \lim \limits_{n \to \infty} n \left[ \frac{1}{(n + 1)^2} + \frac{1}{(n+2)^2} + \cdots + \frac{1}{(n+n)^2} \right] \\
           & = \int_{0}^{1} \left(\frac{1}{1 + x}\right)^2 dx                                                                  \\
           & = - (x+1)^{-1} \bigg|_0^1                                                                                         \\
           & = - \frac{1}{2} - (-1)                                                                                            \\
           & = \frac{1}{2}
        \end{align*}

  \item (3)

        与(2)类似,是函数$f(x) = \left(\frac{1}{1 + x^2}\right)^2$在区间$[0, 1]$上的一个积分和。所以
        \begin{align*}
          \int_{0}^{1} \frac{1}{1 + x^2} dx
           & = arc tan x \bigg|_0^1 \\
           & = \frac{\pi}{4}
        \end{align*}

  \item (4)
        \begin{align*}
            & \lim \limits_{n \to \infty} \frac{1}{n} (sin \frac{\pi}{n} + sin \frac{2\pi}{n} + \cdots + sin \frac{n-1}{n}\pi)                       \\
          = & \lim \limits_{n \to \infty} \frac{1}{n} (sin \frac{0 \pi}{n} + sin \frac{\pi}{n} + sin \frac{2\pi}{n} + \cdots + sin \frac{n-1}{n}\pi) \\
          = & \lim \limits_{n \to \infty} \sum \limits_{i = 1}^{n} \frac{1}{n} \cdot sin \left(\frac{(i - 1)\pi}{n}\right)
        \end{align*}
        这个和式是函数$f(x) = sin x \pi$在区间$[0, 1]$上的一个积分和。所以
        \begin{align*}
           & \lim \limits_{n \to \infty} \frac{1}{n} (sin \frac{\pi}{n} + sin \frac{2\pi}{n} + \cdots + sin \frac{n-1}{n}\pi) \\
           & = \int_{0}^{1} sin x \pi dx                                                                                      \\
           & = - \frac{1}{\pi} cos \pi x \bigg|_0^1                                                                           \\
           & = \frac{1}{\pi} - (- \frac{1}{\pi})                                                                              \\
           & = \frac{2}{\pi}
        \end{align*}
\end{itemize}

\section*{3}

$f$在$[a,b]$上可积，那么对任意$\epsilon > 0$，存在$\delta > 0$，使得对任意分割$T$，
以及在其上任意选取的点集$\{\xi_i\}$，只需满足$||T|| < \delta$时，就有
\begin{align*}
  \left| \sum\limits_{i=1}^n f(\xi_i) \Delta x_i - J \right| < \epsilon
\end{align*}
其中$J = \int_a^b f(x) dx$。

于是，把$F'(x) \neq f(x)$的有限个点作为分割点，并在此基础上找一个分割$T$，并满足$||T|| < \delta$。
此时在每一个小区间上，$F(x)$的拉格朗日中值定理任然成立。

参考定理9.1的证明，容易得到
\begin{align*}
  F(b) - F(a) = \int_a^b f(x) dx
\end{align*}

\end{document}